
\textbf{Définition 3.9.} \uline{Soit} @@M0@@ \uline{un} @@M1@@-\uline{anneau}~; \uline{on appelle} @@M2@@-@@M3@@-algèbre augmentée \uline{la donnée d'une} @@M4@@-@@M5@@-\uline{algèbre} @@M6@@ \uline{et d'un} @@M7@@-\uline{homomorphisme de} @@M8@@-\uline{algèbres} @@M9@@~; un homomorphisme de @@M10@@-@@M11@@-algèbres augmentées \uline{est un} @@M12@@-\uline{homomorphisme compatible aux augmentations}.

\textbf{Exemple 3.9.1.} : Soit @@M13@@ un topos commutativement localement annelé, et @@M14@@ le groupe de Grothendieck de la catégorie des faisceaux de @@M15@@-modules localement libres de type fini sur @@M16@@. Nous avons vu que @@M17@@ peut être muni d'une structure de pré-@@M18@@-anneau (2.2.b, car tout facteur direct d'un libre de type fini est alors localement libre de type fini) et que @@M19@@ est canoniquement muni d'une structure d'anneau binômial. On définit un @@M20@@-homomorphisme de @@M21@@ dans @@M22@@ de la façon suivante : soit @@M23@@ une section du faisceau constant @@M24@@ sur @@M25@@~; elle définit une partition de l'objet final de @@M26@@ en une somme d'ouverts disjoints @@M27@@, la section @@M28@@ ayant pour valeur @@M29@@ sur l'ouvert @@M30@@~; soit alors @@M31@@ le faisceau égal à @@M32@@ sur @@M33@@ pour @@M34@@, et à @@M35@@ sur @@M36@@ pour @@M37@@, et soit @@M38@@ le faisceau égal à @@M39@@ sur @@M40@@ pour @@M41@@, et à @@M42@@ sur @@M43@@ pour @@M44@@~; on associe à @@M45@@ l'élément @@M46@@ de @@M47@@.

Il est immédiat de vérifier qu'on obtient bien ainsi un @@M48@@-homomorphisme de @@M49@@ dans @@M50@@~; on peut alors définir l'augmentation de @@M51@@ dans @@M52@@, en associant à tout faisceau localement libre de type fini son rang, qui est une section de @@M53@@~; l'additivité du rang entraîne qu'on définit ainsi une fonction sur @@M54@@, à valeurs dans @@M55@@. On a donc muni @@M56@@ d'une structure de pré-@@M57@@-algèbre augmentée. C'est toujours de celle-ci qu'il sera question dans la suite.

